\section{Conclusion}
\label{sec:conclusion}

We introduced level-aware hybrid decoding for generative retrieval — a general inference framework that mixes dense and autoregressive scores at each codebook level during beam expansion, motivated by the coarse-to-fine residual structure of RQ-VAE tokenization. The per-level $\alpha$ schedule, whether grid-searched or learned, consistently assigns higher weight to dense guidance at coarser levels, confirming that residual entropy provides a principled basis for the mixing schedule. Our method achieves statistically significant improvements over both vanilla TIGER and LIGER on the majority of datasets tested.

The key insight is that partial RQ reconstructions $r_l^{\text{cand}}$ carry increasing item-identity signal as $l$ grows, making them suitable scoring targets at each level rather than waiting for full generation. This principle — using intermediate reconstruction states as dense scoring anchors — generalizes beyond recommendation to any retrieval task using hierarchical discrete representations.

\paragraph{Limitations.} Level-aware decoding adds dot-product overhead proportional to beam size; for very large codebooks this could become significant. The per-level $\alpha$ schedule requires either a grid search (inference-only, cheap) or a brief fine-tuning pass (one epoch, cheap but adds a training step). LIGER comparison relies on cited paper numbers rather than our own reproduction.

\paragraph{Future work.} Per-level mixing with encoder-side dense models, extension to multi-level RQ-VAEs with $L > 4$, and online adaptation of $\alpha_l$ based on user-specific entropy profiles.
